%==============================================================================
%  MATH 231 -- Linear Algebra I (Queens College, Fall 2026)
%  HOMEWORK 2 -- Norms, Bases, R^n, and the Gradient
%
%  SCOPE.  This file covers five lectures:
%     lectures/vector-unit/MATH231_vectors_2_dot_product_and_projection  (S10-S12)
%     lectures/vector-unit/MATH231_vectors_3_norms_and_metrics           (S13-S15)
%     lectures/vector-unit/MATH231_vectors_4_basis_and_span              (S16-S20)
%     lectures/vector-unit/MATH231_vectors_5_orthogonality_and_subspaces (S21-S24)
%     lectures/vector-calc-unit/MATH231_veccalc_1_derivatives_and_the_gradient
%                                                                  (S1-S7)
%  Nothing here uses veccalc S8 onward.  The draft's three problems on the
%  object-type of the gradient, the directional derivative, and gradient
%  descent / linear approximation were held back for a later assignment; they
%  are in plan/future_homework/ (gitignored).
%
%  WRITTEN ONLY.  There is no programming part.  The NumPy part, with its
%  Python submission rules, is a separate, optional assignment:
%  homework/MATH231_extra_credit_1.tex.  It refers to this assignment's problems
%  by number, so renumbering a problem here means updating it there too.
%
%  PART A.  Unit vectors, unknown components, an application (work and forces
%  on a ramp), the failure of cancellation, the parallelogram and polarization
%  identities, and the algebra of the projection map.
%
%  PART E IS THE R^n PART.  Its applications: document similarity (E.4),
%  mean/variance/correlation as a projection and a cosine (E.5), and the
%  geometry of the high-dimensional cube (E.6).
%
%  NUMBERING.  As in Homework 1: parts are lettered, problems numbered by part.
%  Section references are to the VECTOR UNIT unless a problem is in Part F,
%  whose preamble says that its bare section numbers refer to the vector
%  calculus unit (which restarts at S1).
%
%  DELIBERATE OMISSIONS, so that nothing here outruns the lectures:
%   - No matrices.
%   - No cross product, and no general n x n independence test.  Independence
%     of three vectors in R^3 (C.7) is checked from the definition in the
%     S24.2 table.
%   - Spanning of R^3 by three independent vectors is "announced rather than
%     established" in S21.4.  So D.1(c) makes the student substitute the
%     coordinates back, which is what actually shows the vector is that
%     combination.
%   - No second derivatives, Hessians, or R^n calculus (veccalc Part 2).
%
%  NO ANSWER KEY, by design, as in Homework 1.
%
%  TWO THINGS TO SET BEFORE HANDING THIS OUT:
%     \duedate  -- currently a placeholder.
%     The submission path in the "What to submit" callout, if you want a
%     different directory layout than homework/hw2/ in the student repo.
%
%  Compile with pdflatex or tectonic (standard TeX Live packages only).
%==============================================================================

\documentclass[11pt]{article}

\usepackage[margin=1in]{geometry}
\usepackage{amsmath,amssymb}
\usepackage[dvipsnames]{xcolor}
\usepackage{enumitem}
\usepackage{booktabs}
\usepackage{array}
\usepackage[hidelinks]{hyperref}
\usepackage{titlesec}
\usepackage{needspace}

%------------------------------------------------------------------ style ----
\titleformat{\section}{\large\bfseries\sffamily\color{RoyalBlue}}
                      {Part \thesection.}{0.6em}{}
\titleformat{\subsection}{\normalsize\bfseries\sffamily}{\thesubsection.}{0.5em}{}
\titlespacing*{\section}{0pt}{16pt}{7pt}

\renewcommand{\thesection}{\Alph{section}}

% Keeps long inline formulas from pushing a line into the margin.
\setlength{\emergencystretch}{3em}

% The problem heading is a line of its own, so that part (a) -- or the opening
% sentence -- always starts below it.  (An amsthm theorem cannot promise this: a
% list that opens the body runs its first item onto the heading line.)  The
% extra air above and below keeps consecutive questions reading as separate
% blocks rather than one continuous column of text.
\makeatletter
\newcounter{prob}[section]
\renewcommand{\theprob}{\thesection.\arabic{prob}}
\newenvironment{prob}[1][]
  {\par\addvspace{20pt plus 6pt minus 4pt}%
   \needspace{4\baselineskip}%
   \refstepcounter{prob}%
   \noindent\textbf{Problem~\theprob}%
   \if\relax\detokenize{#1}\relax\else\ (#1)\fi\textbf{.}%
   \par\nobreak\@afterindentfalse\@afterheading}
  {\par\addvspace{20pt plus 6pt minus 4pt}}
\makeatother

\newcommand{\heads}[1]{\par\smallskip\noindent
  {\small\sffamily\textcolor{BrickRed}{\textbf{Watch out.}}\ #1}\par\smallskip}
\newcommand{\hint}[1]{\par\smallskip\noindent
  {\small\sffamily\textcolor{RoyalBlue}{\textbf{Hint.}}\ #1}\par\smallskip}

% shaded call-out box.  The body is collected into a savebox first, because
% \colorbox needs its whole argument at once.
\definecolor{softbg}{HTML}{F2F6FA}
\newsavebox{\qcbox}
\newenvironment{callout}[1]
  {\par\medskip\needspace{5\baselineskip}%
   \begin{lrbox}{\qcbox}%
   \begin{minipage}{\dimexpr\textwidth-2\fboxsep-4pt}%
   \setlength{\parskip}{0.4em}%
   \textbf{\sffamily\color{RoyalBlue}#1}\par}
  {\end{minipage}\end{lrbox}%
   \begin{center}\colorbox{softbg}{\usebox{\qcbox}}\end{center}\medskip}
%--------------------------------------------------------------- notation ----
\newcommand{\R}{\mathbb{R}}
\newcommand{\vc}[1]{\mathbf{#1}}
\newcommand{\uv}[1]{\hat{\vc{#1}}}
\newcommand{\one}{\mathbf{1}}
\newcommand{\norm}[1]{\lVert #1\rVert}
\newcommand{\ip}[2]{\langle #1,#2\rangle}
\newcommand{\proj}[2]{\operatorname{proj}_{#1}(#2)}
\newcommand{\comp}[2]{\operatorname{comp}_{#1}(#2)}
\newcommand{\Sph}[1]{\mathbb{S}^{#1}}
\newcommand{\ivec}{\vec{\imath}}
\newcommand{\jvec}{\vec{\jmath}}
\newcommand{\kvec}{\vec{k}}
\DeclareMathOperator{\spanop}{span}
\newcommand{\grad}{\nabla}
\newcommand{\pd}[2]{\frac{\partial #1}{\partial #2}}
\newcommand{\dd}[2]{\frac{d #1}{d #2}}

%--------------------------------------------------------- fill this in ------
\newcommand{\duedate}{to be announced on the class Discord}

\title{\vspace{-1.2cm}\textbf{\sffamily MATH 231 -- Linear Algebra I}\\[4pt]
       {\large\sffamily Homework 2 \textbullet\ Norms, Bases, $\R^{n}$, and
        the Gradient}\\[2pt]
       {\normalsize\sffamily Kiely Hall 326 \textbullet\
        Instructor: Kevin Bedoya}}
\author{}
\date{}

\begin{document}
\maketitle
\vspace{-1.6cm}

\noindent\textbf{Due:} \duedate.

\medskip
\begin{center}
\renewcommand{\arraystretch}{1.25}
\begin{tabular}{@{}l l >{\raggedright\arraybackslash}p{7.2cm}@{}}
\toprule
\textbf{Lecture} & \textbf{Sections} & \textbf{Topics} \\
\midrule
Vector Unit, Part 2 & \S10--\S12 &
  \emph{revisited}: normalization and unit vectors; the direction--magnitude
  decomposition; projection and the orthogonal piece, used in new ways \\[3pt]
Vector Unit, Part 3 & \S13--\S15 &
  Cauchy--Schwarz and its equality case; the metric axioms; the Manhattan and
  Chebyshev metrics; the $p$-norm and $p\to\infty$; unit balls; the norm
  axioms \textbf{N1}--\textbf{N3} \\[3pt]
Vector Unit, Part 4 & \S16--\S20 &
  the standard basis; linear combinations, span, basis, coordinates;
  dimension; linear independence; subspaces; $\R^{3}$ \\[3pt]
Vector Unit, Part 5 & \S21--\S24 &
  orthogonal vectors are independent; coordinates by projection; making a
  basis orthogonal; the Pythagorean identity; $\R^{n}$ and every definition in
  $\Sigma$ notation \\[3pt]
Vector Calculus Unit, Part 1 & \S1--\S7 &
  the derivative and its rules; $e^{x}$ and $\ln x$ and their derivatives;
  functions of two variables; partial derivatives; the gradient \\
\bottomrule
\end{tabular}
\end{center}

\noindent\textbf{Grading.} Every problem is worth the same number of points.

\medskip
\noindent\textbf{On the proof problems} (A.1(d), A.4(c)--(d), A.5, A.6(a)--(c),
B.3, B.4, B.5, the justifications in C.4, C.5, C.6, D.3, D.4, D.5, E.2, E.3(c),
E.5(a), E.7, F.3, and F.6(c)--(d)). None of them requires a clever idea. What
they require is that you write down what you are assuming, work in components or
from the stated laws, justify each equality by naming the fact that licenses it,
and end with a sentence saying what you have shown. A chain of equalities with
no words is not a proof; neither is a paragraph of words with no computation.
Several of them ask you \emph{not} to use components: that restriction is the
point of the problem, because an argument that uses only the laws works in
$\R^{n}$ unchanged.

\medskip
\noindent\textbf{Collaboration} is encouraged --- work at a board together, argue
about the problems. What you submit must be written by you, in your own words.

\begin{callout}{What to submit}
Push your written work for Parts A--F to your own MATH 231 GitHub repository,
as a single PDF named \texttt{hw2.pdf} in a folder named \texttt{homework/hw2/}.
Handwritten and scanned or photographed is completely fine, as long as it is
legible and the pages are in order.

\textbf{PDF is the only accepted format.} A \texttt{.docx}, a \texttt{.jpg}, a
\texttt{.png}, a \texttt{.heic}, or a link to a file hosted somewhere else will
not be graded. Phone scanning apps export PDF directly, and every word processor
can ``Save as PDF''.

\textbf{Type Parts A--F up instead and you earn $+5$ points of extra credit on
this assignment.} Any typesetting system counts --- \LaTeX, Word's equation
editor, or anything else that produces readable mathematics --- and the extra
credit is for the whole of Parts A--F, not for typing up some of it. The file
you submit must still be a PDF.

If you have not set up the repository yet, Setup Guide 1 (\S\,``Create your
folders'' and \S\,``Save your work and send it to GitHub'') is the walkthrough.
Only pushed work can be graded.
\end{callout}

\begin{callout}{Label every answer}
Label each answer with the number of the problem it answers --- \texttt{A.3(d)},
\texttt{C.2(e)}, \texttt{F.4(h)} --- and put your answers \textbf{in the same
order as they appear in this assignment}. A grader who cannot tell which
question a piece of work belongs to cannot award it credit, and that is the most
common way a correct answer loses points.

\textbf{You do not have to copy the question out.} This applies to handwritten
and typed submissions alike: the label alone is enough. Write \texttt{A.3(d)},
then your work. Do not reproduce the wording of the problem.
\end{callout}

%==============================================================================
\section{Dot product and projection, revisited}
%==============================================================================
\emph{Vector Unit, Part 2.} For $\vc{u}\neq\vc{0}$,
\[
  \uv{u} \;=\; \frac{\vc{u}}{\norm{\vc{u}}},
  \qquad
  \comp{\vc{u}}{\vc{v}} \;=\; \frac{\ip{\vc{u}}{\vc{v}}}{\norm{\vc{u}}},
  \qquad
  \proj{\vc{u}}{\vc{v}} \;=\; \frac{\ip{\vc{u}}{\vc{v}}}{\norm{\vc{u}}^{2}}\,\vc{u} ,
\]
and the laws of the dot product are \textbf{D1} symmetry, \textbf{D2}
additivity, \textbf{D3} homogeneity, \textbf{D4} positive definiteness,
introduced in \S11.8 (``The algebraic laws of the dot product'') of the Vector
Unit, Part 2 lecture, \emph{Orientation, the Dot Product, and Projection}.

%------------------------------------------------------------------------------
\begin{prob}[Unit vectors and the direction--magnitude decomposition]
\begin{enumerate}[label=\textbf{(\alph*)},leftmargin=2.6em,itemsep=7pt,topsep=6pt]
  \item Normalize each vector below and check that the result has norm $1$.
        Then write each vector in the direction--magnitude form
        $\vc{u}=\norm{\vc{u}}\,\uv{u}$ of \S12.2:
        \[
          [\,-5,\,12\,],\qquad [\,2,\,-2\,],\qquad [\,0,\,-7\,],\qquad
          [\,\sqrt3,\,1\,].
        \]
  \item Write down the vector of length $10$ pointing the same way as
        $[\,3,\,-4\,]$, and the vector of length $3$ pointing \emph{opposite}
        to $[\,1,\,2\,]$. Build each one as (scalar)$\times$(unit vector).
  \item Recall from \S12.1 that $\Sph{1}$ is the \textbf{unit circle}, the set of
        all unit vectors in the plane:
        \[
          \Sph{1} \;=\; \bigl\{\,\vc{u}\in\R^{2} \;:\; \norm{\vc{u}}=1\,\bigr\}.
        \]
        Which of $[\,0.6,\,-0.8\,]$, $[\,\tfrac12,\,\tfrac12\,]$,
        $[\,-1,\,0\,]$ and $[\,\cos2,\,\sin2\,]$ lie in $\Sph{1}$? Justify each.
  \item Let $\vc{u}\neq\vc{0}$ and $\alpha\neq0$. Show that normalizing
        $\alpha\vc{u}$ gives $\uv{u}$ when $\alpha>0$ and $-\uv{u}$ when
        $\alpha<0$. Which rule from \S8 does the work?
\end{enumerate}
\end{prob}

%------------------------------------------------------------------------------
\begin{prob}[Solving for an unknown component]
Find every real number $k$ satisfying each condition. Show the equation you
solve.
\begin{enumerate}[label=\textbf{(\alph*)},leftmargin=2.6em,itemsep=7pt,topsep=6pt]
  \item $[\,k,\,3\,]$ is orthogonal to $[\,2,\,-4\,]$.
  \item $[\,k,\,1\,]$ is orthogonal to $[\,k,\,-4\,]$.
  \item The angle between $[\,k,\,1\,]$ and $[\,1,\,0\,]$ is $\pi/3$.
  \item $\cos\bigl([\,1,\,k\,],[\,3,\,3\,]\bigr)=\tfrac12$.
\end{enumerate}

\heads{In (c) and (d) you will want to square both sides. Squaring can create a
``solution'' of the squared equation that fails the original one. Check every
candidate in the original equation and say why you discard any that fail.}
\end{prob}

%------------------------------------------------------------------------------
\begin{prob}[Work, and a box on a ramp]
In physics a force is a \emph{vector} $\vc{F}\in\R^{2}$: its direction is the
way it pushes, and its norm $\norm{\vc{F}}$ is how hard. Call a force
\emph{constant} when it is the same vector at every moment of the motion. If a
constant force acts on an object that moves along a straight line, from its
starting point to its finishing point, the \emph{displacement} is the vector
$\vc{s}\in\R^{2}$ from start to finish, and the \textbf{work} done by the force
is the scalar
\[
  W \;=\; \ip{\vc{F}}{\vc{s}} ,
\]
a single real number --- in joules, when $\vc{F}$ is in newtons and $\vc{s}$ in
metres. So in this problem $\vc{F}$, $\vc{s}$, $\vc{d}$, $\vc{G}$ and $\vc{P}$
are vectors in $\R^{2}$, while $W$, every norm, and every component such as
$\comp{\vc{d}}{\vc{G}}$ are scalars. A box sits on a ramp that rises in the
direction $\vc{d}=[\,4,\,3\,]$, and gravity pulls on it with the constant force
$\vc{G}=[\,0,\,-50\,]$.

\begin{enumerate}[label=\textbf{(\alph*)},leftmargin=2.6em,itemsep=7pt,topsep=6pt]
  \item Compute $\comp{\vc{d}}{\vc{G}}$, $\proj{\vc{d}}{\vc{G}}$, and the
        orthogonal piece $\vc{G}-\proj{\vc{d}}{\vc{G}}$. Check that the
        orthogonal piece is orthogonal to $\vc{d}$.
  \item Compute the norms of the two pieces. One of them pulls the box down the
        ramp and the other presses it into the ramp surface. Which is which?
  \item The box is pushed from $(0,0)$ to $(8,6)$. Compute the work done by
        gravity, and explain what its sign means.
  \item Over the same displacement a worker pushes with $\vc{P}=[\,30,\,10\,]$.
        Compute the work two ways: as $\ip{\vc{P}}{\vc{s}}$, and as
        $\comp{\vc{s}}{\vc{P}}\cdot\norm{\vc{s}}$. Explain why the second
        formula says that only the part of $\vc{P}$ along $\vc{s}$ does work.
  \item Of all forces of magnitude $40$, which one does the most work over this
        displacement, and how much? Name a nonzero force that does no work at
        all. Justify both with $\ip{\vc{F}}{\vc{s}}=\norm{\vc{F}}\norm{\vc{s}}\cos\theta$.
\end{enumerate}
\end{prob}

%------------------------------------------------------------------------------
\begin{prob}[No cancellation, and what an equation of dot products says]
Let $\vc{u}=[\,2,\,-1\,]$.
\begin{enumerate}[label=\textbf{(\alph*)},leftmargin=2.6em,itemsep=7pt,topsep=6pt]
  \item Find three different vectors $\vc{v}$ with $\ip{\vc{u}}{\vc{v}}=5$.
  \item Compute $\proj{\vc{u}}{\vc{v}}$ for each of your three. What do you
        notice? Describe the set of \emph{all} $\vc{v}$ with
        $\ip{\vc{u}}{\vc{v}}=5$ geometrically.
  \item Prove that for all $\vc{u},\vc{v},\vc{w}\in\R^{2}$: if
        $\ip{\vc{u}}{\vc{v}}=\ip{\vc{u}}{\vc{w}}$, then $\vc{u}$ is
        orthogonal to $\vc{v}-\vc{w}$. Use only \textbf{D1}--\textbf{D4}.
  \item Prove: if $\ip{\vc{x}}{\vc{v}}=\ip{\vc{x}}{\vc{w}}$ for
        \emph{every} $\vc{x}\in\R^{2}$, then $\vc{v}=\vc{w}$.
        \hint{You are allowed to choose $\vc{x}$. One choice finishes the proof in
        two lines using \textbf{D4}.}
\end{enumerate}
\end{prob}

%------------------------------------------------------------------------------
\begin{prob}[The parallelogram law, polarization, and two theorems from Euclid]
In (a)--(d), work from \textbf{D1}--\textbf{D4} and
$\ip{\vc{w}}{\vc{w}}=\norm{\vc{w}}^{2}$ only --- no components --- so that each
argument holds in $\R^{n}$ as written.
\begin{enumerate}[label=\textbf{(\alph*)},leftmargin=2.6em,itemsep=7pt,topsep=6pt]
  \item \textbf{Parallelogram law.} Prove
        $\norm{\vc{u}+\vc{v}}^{2}+\norm{\vc{u}-\vc{v}}^{2}
         =2\norm{\vc{u}}^{2}+2\norm{\vc{v}}^{2}$.
  \item \textbf{Polarization.} Prove
        $\ip{\vc{u}}{\vc{v}}
         =\tfrac14\bigl(\norm{\vc{u}+\vc{v}}^{2}-\norm{\vc{u}-\vc{v}}^{2}\bigr)$.
  \item Prove
        $\ip{\vc{u}+\vc{v}}{\vc{u}-\vc{v}}=\norm{\vc{u}}^{2}-\norm{\vc{v}}^{2}$,
        and deduce that \emph{the diagonals of a rhombus are orthogonal}. (A
        rhombus is a parallelogram with four equal sides, which is the case
        $\norm{\vc{u}}=\norm{\vc{v}}$.)
  \item \textbf{Thales' theorem.} Let $\vc{a}\neq\vc{0}$ and let $\vc{p}$ be any
        vector with $\norm{\vc{p}}=\norm{\vc{a}}$ and $\vc{p}\neq\pm\vc{a}$. So
        $-\vc{a}$, $\vc{a}$ and $\vc{p}$ are three points on the circle of
        radius $\norm{\vc{a}}$ centred at the origin, and $-\vc{a}$, $\vc{a}$
        are the ends of a diameter. Use (c) to prove that
        $\vc{p}-\vc{a}$ is orthogonal to $\vc{p}+\vc{a}$. Then say which angle
        of which triangle has been shown to be a right angle.
  \item Check (a) and (b) numerically on $\vc{u}=[\,3,\,1\,]$ and
        $\vc{v}=[\,-1,\,2\,]$, and check (d) on $\vc{a}=[\,5,\,0\,]$,
        $\vc{p}=[\,3,\,4\,]$.
\end{enumerate}
\end{prob}

%------------------------------------------------------------------------------
\begin{prob}[Three facts about projection]
Let $\vc{u}\neq\vc{0}$. Prove each statement for all $\vc{v},\vc{w}\in\R^{2}$ and
scalars $\alpha,\beta$, using the formula for $\proj{\vc{u}}{\vc{v}}$ and
\textbf{D1}--\textbf{D4}.
\begin{enumerate}[label=\textbf{(\alph*)},leftmargin=2.6em,itemsep=7pt,topsep=6pt]
  \item \textbf{Only the line matters.} For every $\alpha\neq0$,
        $\proj{\alpha\vc{u}}{\vc{v}}=\proj{\vc{u}}{\vc{v}}$. But
        $\comp{\alpha\vc{u}}{\vc{v}}$ equals $\comp{\vc{u}}{\vc{v}}$ when
        $\alpha>0$ and $-\comp{\vc{u}}{\vc{v}}$ when $\alpha<0$.
  \item \textbf{A shadow of a shadow.}
        $\proj{\vc{u}}{\proj{\vc{u}}{\vc{v}}}=\proj{\vc{u}}{\vc{v}}$. Say what
        this means geometrically.
  \item \textbf{Projection is linear in what is projected.}
        $\proj{\vc{u}}{\vc{v}+\vc{w}}=\proj{\vc{u}}{\vc{v}}+\proj{\vc{u}}{\vc{w}}$
        and $\proj{\vc{u}}{\beta\vc{v}}=\beta\,\proj{\vc{u}}{\vc{v}}$.
  \item \textbf{But not in the direction projected onto.} Show by computing
        that $\proj{\vc{u}+\vc{w}}{\vc{v}}\neq\proj{\vc{u}}{\vc{v}}+\proj{\vc{w}}{\vc{v}}$
        for $\vc{u}=[\,1,\,0\,]$, $\vc{w}=[\,1,\,1\,]$, $\vc{v}=[\,0,\,1\,]$.
\end{enumerate}
\end{prob}

%==============================================================================
\section{Cauchy--Schwarz, metrics, and the family of norms}
%==============================================================================
\emph{Vector Unit, Part 3.} For $\vc{v}=[\,v_1,\,v_2\,]$ and $p\ge1$,
\[
  \norm{\vc{v}}_1=\lvert v_1\rvert+\lvert v_2\rvert,
  \qquad
  \norm{\vc{v}}_2=\sqrt{v_1^{2}+v_2^{2}},
  \qquad
  \norm{\vc{v}}_\infty=\max\bigl\{\lvert v_1\rvert,\lvert v_2\rvert\bigr\},
  \qquad
  \norm{\vc{v}}_p=\bigl(\lvert v_1\rvert^{p}+\lvert v_2\rvert^{p}\bigr)^{1/p},
\]
and each norm gives a metric $d_p(\vc{x},\vc{y})=\norm{\vc{x}-\vc{y}}_p$. A bare
$\norm{\vc{v}}$ still means $\norm{\vc{v}}_2$. The metric axioms are
\textbf{M1}--\textbf{M4} (\S14.1), the norm axioms \textbf{N1}--\textbf{N3}
(\S15.1). Leave answers exact unless a decimal is asked for.

%------------------------------------------------------------------------------
\begin{prob}[The family of norms, by hand]
Let $\vc{v}$ range over the four vectors
\[
  \text{(a) }[\,-3,\,4\,],\qquad
  \text{(b) }[\,5,\,-12\,],\qquad
  \text{(c) }[\,2,\,2\,],\qquad
  \text{(d) }[\,-7,\,0\,],
\]
\begin{enumerate}[label=\textbf{(\roman*)},leftmargin=2.6em,itemsep=7pt,topsep=6pt]
  \item Compute $\norm{\vc{v}}_1$, $\norm{\vc{v}}_2$ and $\norm{\vc{v}}_\infty$,
        and confirm the ordering
        $\norm{\vc{v}}_\infty\le\norm{\vc{v}}_2\le\norm{\vc{v}}_1$ of \S14.8.
  \item For one of the four, all three norms are equal. Which one, and what
        feature of the vector forces it?
  \item For vector (a), compute $\norm{\vc{v}}_p$ for $p=3$, $4$ and $10$, to four
        decimal places. The values should be decreasing toward
        $\norm{\vc{v}}_\infty$. Compute the upper bound $2^{1/p}M$ from the
        squeeze argument of \S14.7 at $p=10$, where
        $M=\norm{\vc{v}}_\infty$, and confirm that $\norm{\vc{v}}_{10}$ lies
        between $M$ and that bound.
\end{enumerate}
\end{prob}

%------------------------------------------------------------------------------
\begin{prob}[Which customer is nearest? It depends on the metric]
A depot sits at $P=(1,-1)$, and three customers at $A=(7,-1)$, $B=(6,1)$ and
$C=(5,3)$.
\begin{enumerate}[label=\textbf{(\alph*)},leftmargin=2.6em,itemsep=7pt,topsep=6pt]
  \item Make a table of $d_1$, $d_2$ and $d_\infty$ from $P$ to each customer.
  \item Three vehicles leave the depot: a \emph{taxi} that must follow a
        rectangular street grid, a \emph{drone} that flies in a straight line,
        and an \emph{overhead crane} whose east--west and north--south motors run
        at the same time and the same speed. Which metric measures the trip for
        each vehicle? Which customer is nearest for each?
  \item On one set of axes, sketch the three ``circles'' of radius $6$ centred at
        $P$ --- the sets $\{\vc{x}: d(P,\vc{x})=6\}$ for $d=d_1,d_2,d_\infty$ ---
        and plot $A$, $B$, $C$.
\end{enumerate}
\end{prob}

%------------------------------------------------------------------------------
\begin{prob}[The Manhattan and Chebyshev metrics are metrics]
A function $d$ that takes two vectors of $\R^{2}$ and returns a real number is a
\textbf{metric} on $\R^{2}$ when it satisfies the metric axioms of \S14.1. The
three you need here are, for all $\vc{x},\vc{y},\vc{z}\in\R^{2}$:
\begin{center}
\renewcommand{\arraystretch}{1.4}
\begin{tabular}{@{}l l l@{}}
\textbf{M1} & $d(\vc{x},\vc{y})=0$ if and only if $\vc{x}=\vc{y}$
            & zero distance only from a point to itself \\
\textbf{M2} & $d(\vc{x},\vc{y})=d(\vc{y},\vc{x})$ & symmetry \\
\textbf{M3} & $d(\vc{x},\vc{y})+d(\vc{y},\vc{z})\ge d(\vc{x},\vc{z})$
            & the triangle inequality \\
\end{tabular}
\end{center}
\S14.3 left the first of the two metrics below as an exercise. Prove both
directly from the definitions and the scalar triangle inequality
$\lvert a+b\rvert\le\lvert a\rvert+\lvert b\rvert$ for real numbers ---
not by citing \S15.2.
\begin{enumerate}[label=\textbf{(\alph*)},leftmargin=2.6em,itemsep=7pt,topsep=6pt]
  \item Prove that $d_1$ satisfies \textbf{M1}, \textbf{M2} and \textbf{M3} on
        $\R^{2}$.
  \item Prove the same for $d_\infty$.
        \hint{For \textbf{M3}, fix a coordinate $i$ and show
        $\lvert x_i-z_i\rvert\le d_\infty(\vc{x},\vc{y})+d_\infty(\vc{y},\vc{z})$.
        Since this holds for $i=1$ and for $i=2$, it holds for the larger of
        the two.}
\end{enumerate}
\end{prob}

%------------------------------------------------------------------------------
\begin{prob}[When Cauchy--Schwarz is an equality]
\S13 showed, using the angle $\theta$, that equality in
$\lvert\ip{\vc{u}}{\vc{v}}\rvert\le\norm{\vc{u}}\norm{\vc{v}}$ happens exactly
for parallel or antiparallel vectors, and the remark ``Cauchy--Schwarz, one more
time'' in \S23.4 settled it again using the Pythagorean identity. Prove it a
third way, from \textbf{D1}--\textbf{D4} alone --- no angle, and no appeal to
Pythagoras --- so that nothing but the laws is needed.
\begin{enumerate}[label=\textbf{(\alph*)},leftmargin=2.6em,itemsep=7pt,topsep=6pt]
  \item Show that if $\vc{v}=\alpha\vc{u}$ for some scalar $\alpha$, then
        $\lvert\ip{\vc{u}}{\vc{v}}\rvert=\norm{\vc{u}}\norm{\vc{v}}$.
  \item Now suppose $\vc{u}\neq\vc{0}$ and
        $\lvert\ip{\vc{u}}{\vc{v}}\rvert=\norm{\vc{u}}\norm{\vc{v}}$. Put
        $c=\ip{\vc{u}}{\vc{v}}/\norm{\vc{u}}^{2}$. Expand
        $\norm{\vc{v}-c\vc{u}}^{2}$ using \textbf{D1}--\textbf{D4} and show
        that
        \[
          \norm{\vc{v}-c\vc{u}}^{2}
          \;=\;
          \norm{\vc{v}}^{2}-\frac{\ip{\vc{u}}{\vc{v}}^{2}}{\norm{\vc{u}}^{2}} .
        \]
        Conclude that $\vc{v}=c\vc{u}$.
\end{enumerate}
\end{prob}

%------------------------------------------------------------------------------
\begin{prob}[The reverse triangle inequality]
Let $\norm{\cdot}$ be \emph{any} norm on $\R^{2}$ --- you may use only
\textbf{N1}--\textbf{N3}. Prove that for all $\vc{u},\vc{v}$,
\[
  \bigl\lvert\,\norm{\vc{u}}-\norm{\vc{v}}\,\bigr\rvert
  \;\le\;
  \norm{\vc{u}-\vc{v}} .
\]
\hint{Write $\vc{u}=(\vc{u}-\vc{v})+\vc{v}$ and apply \textbf{N3} to get one
inequality. Swap the roles of $\vc{u}$ and $\vc{v}$ for the other, and use
\textbf{N2} with $\alpha=-1$ to see that $\norm{\vc{v}-\vc{u}}=\norm{\vc{u}-\vc{v}}$.}
\end{prob}

%==============================================================================
\section{Bases, span, independence, and subspaces}
%==============================================================================
\emph{Vector Unit, Part 4.} Recall that $\alpha\vc{u}+\beta\vc{v}$ is a
\textbf{linear combination}, the \textbf{span} is the set of all of them, and a
\textbf{basis} spans without redundancy (\S16.3). Two vectors are
\textbf{linearly independent} when $\alpha\vc{u}+\beta\vc{v}=\vc{0}$ forces
$\alpha=\beta=0$ (\S18.4); for more vectors the definition is the same with more
terms (the table in \S24.2). For a pair in $\R^{2}$, the number
$D=u_1v_2-u_2v_1$ decides everything: when $D\neq0$ the coordinates of
$\vc{w}$ are
\[
  \alpha=\frac{w_1v_2-w_2v_1}{D},\qquad
  \beta=\frac{u_1w_2-u_2w_1}{D}\qquad(\S18.5).
\]
\S18.5 proves that an independent pair has $D\neq0$. The converse holds too, and
you may use it: if $D\neq0$, apply the elimination of \S18.5 with
$\vc{w}=\vc{0}$, which gives $\alpha D=0$ and $\beta D=0$, so the only
combination producing $\vc{0}$ is the trivial one.
A \textbf{subspace} is a subset that is non-empty, closed under addition, and
closed under scalar multiplication (\S19.3).

\heads{The lecture labels the three subspace conditions S1--S3. Those are
\emph{not} the scalar-multiplication laws S1--S4 of \S8.4. To avoid the clash,
this assignment refers to them by name: non-empty, closed under addition, closed
under scaling.}

%------------------------------------------------------------------------------
\begin{prob}[Coordinates against a basis that is not the standard one]
Let $\vc{u}=[\,2,\,1\,]$ and $\vc{v}=[\,1,\,3\,]$.
\begin{enumerate}[label=\textbf{(\alph*)},leftmargin=2.6em,itemsep=7pt,topsep=6pt]
  \item Compute $D$ and explain why $\{\vc{u},\vc{v}\}$ is a basis of $\R^{2}$.
  \item Find the coordinates of $\vc{w}=[\,5,\,5\,]$ and of
        $\vc{w}=[\,1,\,-2\,]$ by setting up the linear system and eliminating
        by hand, as in \S18.1.
\end{enumerate}
\end{prob}

%------------------------------------------------------------------------------
\begin{prob}[Independent or dependent?]
For each pair, decide whether it is linearly independent. Argue from the
definition --- set up $\alpha\vc{u}+\beta\vc{v}=\vc{0}$ and solve --- and
confirm with $D$. For each \emph{dependent} pair, exhibit scalars $\alpha,\beta$,
not both zero, with $\alpha\vc{u}+\beta\vc{v}=\vc{0}$; then describe the span as
a set, give a basis for it, and state its dimension.
\begin{center}
\renewcommand{\arraystretch}{1.25}
\begin{tabular}{@{}c c c @{\qquad} c c c@{}}
\toprule
 & $\vc{u}$ & $\vc{v}$ & & $\vc{u}$ & $\vc{v}$ \\
\midrule
(a) & $[\,3,\,-6\,]$ & $[\,-1,\,2\,]$ & (d) & $[\,\tfrac12,\,\tfrac34\,]$ & $[\,2,\,3\,]$ \\
(b) & $[\,2,\,5\,]$  & $[\,5,\,2\,]$  & (e) & $[\,\sqrt2,\,1\,]$ & $[\,2,\,\sqrt2\,]$ \\
(c) & $[\,0,\,0\,]$  & $[\,4,\,1\,]$  & (f) & $[\,3,\,1\,]$ & $[\,-1,\,3\,]$ \\
\bottomrule
\end{tabular}
\end{center}
\begin{enumerate}[label=\textbf{(\alph*)},leftmargin=2.6em,itemsep=7pt,topsep=6pt]
  \setcounter{enumi}{6}
  \item Find every $k$ for which $[\,k,\,2\,]$ and $[\,8,\,k\,]$ are
        dependent, and for each such $k$ give the equation of the line they
        span.
\end{enumerate}
\end{prob}

%------------------------------------------------------------------------------
\begin{prob}[A spanning set that is not a basis of $\R^{2}$]
Let $S=\bigl\{[\,1,\,2\,],\ [\,3,\,1\,],\ [\,4,\,3\,]\bigr\}$, a set of three
vectors in $\R^{2}$. It spans $\R^{2}$, but it is not a basis of $\R^{2}$.
\begin{enumerate}[label=\textbf{(\alph*)},leftmargin=2.6em,itemsep=7pt,topsep=6pt]
  \item Show that $[\,4,\,3\,]=[\,1,\,2\,]+[\,3,\,1\,]$, and use it to write
        a combination $a[1,2]+b[3,1]+c[4,3]=\vc{0}$ with $a,b,c$ not all zero.
  \item Write $\vc{w}=[\,7,\,4\,]$ as a combination of the three vectors in two
        different ways.
  \item Show that for \emph{every} real number $t$,
        \[
          (1-t)[\,1,\,2\,]+(2-t)[\,3,\,1\,]+t[\,4,\,3\,]=[\,7,\,4\,] .
        \]
        So $\vc{w}$ has infinitely many sets of ``coordinates'' against $S$.
  \item Which of the three two-element subsets of $S$ are bases of $\R^{2}$?
  \item In two sentences: $S$ spans $\R^{2}$, so why is it not a basis? Which
        half of ``basis'' --- existence of coordinates, or their uniqueness ---
        does it fail?
\end{enumerate}
\end{prob}

%------------------------------------------------------------------------------
\begin{prob}[Subspace or not?]
For each subset of $\R^{2}$, decide whether it is a subspace of $\R^{2}$. If it
is, prove the
three conditions, give a basis, and state its dimension. If it is not, name a
condition that fails and give explicit vectors (and a scalar, if needed) that
show it.
\begin{enumerate}[label=\textbf{(\alph*)},leftmargin=2.6em,itemsep=7pt,topsep=6pt]
  \item $\{(x,y) : 3x-5y=0\}$
  \item $\{(x,y) : x+y=1\}$
  \item $\{(x,y) : xy=0\}$
  \item $\{(x,y) : y=x^{2}\}$
  \item $\{(x,y) : x \text{ and } y \text{ are both integers}\}$
  \item $\{(x,y) : 2x=y \text{ and } x=3y\}$
  \item $\{(x,y) : x^{2}+y^{2}\le 1\}$
\end{enumerate}
\end{prob}

%------------------------------------------------------------------------------
\begin{prob}[Three short proofs about span and independence]
Work from the definitions. Do not write vectors in components.
\begin{enumerate}[label=\textbf{(\alph*)},leftmargin=2.6em,itemsep=7pt,topsep=6pt]
  \item Prove: if $\vc{w}\in\spanop\{\vc{u},\vc{v}\}$, then
        $\spanop\{\vc{u},\vc{v},\vc{w}\}=\spanop\{\vc{u},\vc{v}\}$.
        Prove both inclusions. (This is the general form of \S18.6, where
        dropping a redundant vector cost nothing.)
  \item Prove: if $\vc{u}$ and $\vc{v}$ are linearly independent, then so are
        $\vc{u}+\vc{v}$ and $\vc{u}-\vc{v}$.
  \item Prove the converse of (b): if $\vc{u}+\vc{v}$ and $\vc{u}-\vc{v}$ are
        linearly independent, then so are $\vc{u}$ and $\vc{v}$.
        \hint{Write $\vc{a}=\vc{u}+\vc{v}$ and $\vc{b}=\vc{u}-\vc{v}$, express
        $\vc{u}$ and $\vc{v}$ in terms of $\vc{a}$ and $\vc{b}$, and substitute
        into $\alpha\vc{u}+\beta\vc{v}=\vc{0}$.}
\end{enumerate}
\end{prob}

%------------------------------------------------------------------------------
\begin{prob}[Intersections and unions of subspaces]
Let $S$ and $T$ be subspaces of $\R^{2}$.
\begin{enumerate}[label=\textbf{(\alph*)},leftmargin=2.6em,itemsep=7pt,topsep=6pt]
  \item Prove that the intersection $S\cap T$ is a subspace.
  \item Prove: if $S\cup T$ is a subspace, then $S\subseteq T$ or
        $T\subseteq S$.
        \hint{Suppose neither holds, and pick $\vc{s}\in S$ with $\vc{s}\notin T$
        and $\vc{t}\in T$ with $\vc{t}\notin S$. Then $\vc{s}+\vc{t}$ lies in
        $S\cup T$, so it lies in $S$ or in $T$. Show that either case leads to a
        contradiction.}
\end{enumerate}
\end{prob}

%------------------------------------------------------------------------------
\begin{prob}[Spans in $\R^{3}$]
Let $\vc{a}=[\,1,\,1,\,0\,]$ and $\vc{b}=[\,0,\,1,\,1\,]$.
\begin{enumerate}[label=\textbf{(\alph*)},leftmargin=2.6em,itemsep=7pt,topsep=6pt]
  \item Show that $\vc{a}$ and $\vc{b}$ are linearly independent.
  \item Decide whether each of $[\,1,\,2,\,1\,]$, $[\,1,\,0,\,0\,]$ and
        $[\,2,\,-1,\,-3\,]$ lies in $\spanop\{\vc{a},\vc{b}\}$. For each that
        does, give the coefficients. For each that does not, show that the
        system has no solution.
  \item Show that $\spanop\{\vc{a},\vc{b}\}=\{[\,x,\,y,\,z\,] : y=x+z\}$. What
        kind of geometric object is it, and what is its dimension?
  \item Let $\vc{c}=[\,1,\,0,\,0\,]$. Using the definition in the table of
        \S24.2, show that $\vc{a}$, $\vc{b}$, $\vc{c}$ are linearly
        independent, and write $[\,3,\,5,\,4\,]$ as a combination of them.
  \item Is the plane $z=1$ a subspace of $\R^{3}$? One sentence.
\end{enumerate}
\end{prob}

%==============================================================================
\section{Orthogonality and the Pythagorean identity}
%==============================================================================
\emph{Vector Unit, Part 5.} Nonzero, mutually orthogonal vectors are linearly
independent (\S21.2). Against an orthogonal basis, coordinates need no system:
the coefficient of each basis vector $\vc{b}$ is
$\ip{\vc{b}}{\vc{w}}/\norm{\vc{b}}^{2}$, and $\vc{w}$ is the sum of its shadows on
the basis vectors (\S21.6). And if $\vc{u}\perp\vc{v}$ then
$\norm{\vc{u}+\vc{v}}^{2}=\norm{\vc{u}}^{2}+\norm{\vc{v}}^{2}$ (\S23).

\medskip
\noindent In $\R^{3}$, \S21.4 proves that three mutually orthogonal nonzero
vectors are independent, but only \emph{announces} that they span. So wherever a
problem below asks for coordinates in $\R^{3}$, substituting them back is part of
the answer: that step is what shows the vector really is that combination.

%------------------------------------------------------------------------------
\begin{prob}[Coordinates by projection alone]
\begin{enumerate}[label=\textbf{(\alph*)},leftmargin=2.6em,itemsep=7pt,topsep=6pt]
  \item Let $\vc{u}=[\,3,\,4\,]$ and $\vc{v}=[\,-8,\,6\,]$. Check that they are
        nonzero and orthogonal, and say which result makes them a basis of
        $\R^{2}$. Then find the coordinates of $[\,2,\,11\,]$, $[\,-1,\,7\,]$
        and $[\,5,\,0\,]$ with one dot product and one division per coordinate,
        and check each by substituting back.
  \item For $[\,5,\,0\,]$, find the coordinates again with the formula of
        \S18.5. You should get the same numbers with noticeably more work.
  \item Let $\vc{u}=[\,1,\,1,\,1\,]$, $\vc{v}=[\,1,\,-1,\,0\,]$ and
        $\vc{w}=[\,1,\,1,\,-2\,]$. Check that they are mutually orthogonal and
        nonzero. Find the coefficients of $\vc{x}=[\,3,\,5,\,-2\,]$ and of
        $\vc{y}=[\,0,\,0,\,6\,]$ by projection, and substitute back.
  \item Write $\vc{x}$ from (c) as the sum of its three shadows
        $\proj{\vc{u}}{\vc{x}}+\proj{\vc{v}}{\vc{x}}+\proj{\vc{w}}{\vc{x}}$,
        listing each shadow as a vector.
\end{enumerate}
\end{prob}

%------------------------------------------------------------------------------
\begin{prob}[Manufacturing an orthogonal basis]
The \emph{lookahead} at the end of \S21.6 promised that any basis can be
converted into an orthogonal one by subtracting projections. Here is the first
step of that procedure.
\begin{enumerate}[label=\textbf{(\alph*)},leftmargin=2.6em,itemsep=7pt,topsep=6pt]
  \item Let $\vc{u}=[\,3,\,1\,]$ and $\vc{v}=[\,2,\,4\,]$. Show that they are
        independent but not orthogonal. Compute
        $\vc{v}'=\vc{v}-\proj{\vc{u}}{\vc{v}}$ and check that
        $\vc{v}'\perp\vc{u}$.
  \item Find the coordinates of $\vc{w}=[\,7,\,9\,]$ against the orthogonal
        basis $\{\vc{u},\vc{v}'\}$ by projection. Then use
        $\vc{v}'=\vc{v}-\proj{\vc{u}}{\vc{v}}$ to rewrite the result as a
        combination of the original $\vc{u}$ and $\vc{v}$, and check it.
  \item Explain why $\spanop\{\vc{u},\vc{v}'\}=\spanop\{\vc{u},\vc{v}\}$:
        show that each of $\vc{v}'$ and $\vc{v}$ is a combination of the other
        pair.
\end{enumerate}
\end{prob}

%------------------------------------------------------------------------------
\begin{prob}[Projection and the squared cosine]
Prove that for $\vc{u},\vc{v}\neq\vc{0}$,
\[
  \frac{\norm{\proj{\vc{u}}{\vc{v}}}^{2}}{\norm{\vc{v}}^{2}}
  \;=\;\cos^{2}(\vc{u},\vc{v}) ,
\]
and evaluate both sides for $\vc{u}=[\,1,\,2\,]$ and $\vc{v}=[\,4,\,7\,]$.
Read the left side as ``the fraction of $\vc{v}$'s squared length explained by
the direction $\vc{u}$''. For this pair, is it close to $1$ or to $0$, and what
does that say about the angle?
\end{prob}

%------------------------------------------------------------------------------
\begin{prob}[Pythagoras for three vectors, and for $n$]
Work from the laws alone (\S23.2 does), so that the result holds in any
$\R^{n}$.
\begin{enumerate}[label=\textbf{(\alph*)},leftmargin=2.6em,itemsep=7pt,topsep=6pt]
  \item Prove: if $\vc{u},\vc{v},\vc{w}$ are mutually orthogonal, then
        $\norm{\vc{u}+\vc{v}+\vc{w}}^{2}
         =\norm{\vc{u}}^{2}+\norm{\vc{v}}^{2}+\norm{\vc{w}}^{2}$.
  \item Deduce: if $\vc{u},\vc{v},\vc{w}$ are orthonormal, then
        $\norm{\alpha\vc{u}+\beta\vc{v}+\gamma\vc{w}}^{2}
         =\alpha^{2}+\beta^{2}+\gamma^{2}$. In words: against an orthonormal
        basis, the squared length of a vector is the sum of the squares of its
        coordinates.
  \item State, in $\Sigma$ notation and without proof, the version of (a) for
        $k$ mutually orthogonal vectors $\vc{u}_1,\dots,\vc{u}_k$.
\end{enumerate}
\end{prob}

%------------------------------------------------------------------------------
\begin{prob}[Everything orthogonal to a vector]
\begin{enumerate}[label=\textbf{(\alph*)},leftmargin=2.6em,itemsep=7pt,topsep=6pt]
  \item Prove: if $\vc{w}\perp\vc{u}$ and $\vc{w}\perp\vc{v}$, then $\vc{w}$ is
        orthogonal to every vector in $\spanop\{\vc{u},\vc{v}\}$.
  \item For a vector $\vc{u}\in\R^{2}$, let
        $\vc{u}^{\perp}=\{\vc{x}\in\R^{2} : \ip{\vc{u}}{\vc{x}}=0\}$. Prove that
        $\vc{u}^{\perp}$ is a subspace of $\R^{2}$.
\end{enumerate}
\end{prob}

%==============================================================================
\section{Computing in \texorpdfstring{$\R^{n}$}{Rn}}
%==============================================================================
\emph{Vector Unit, \S20.3 and \S24.} Every definition of the unit holds in
$\R^{n}$ with a $\Sigma$ in it (the table in \S24.2). The ones used most below:
\[
  \ip{\vc{u}}{\vc{v}}=\sum_{i=1}^{n}u_iv_i,
  \qquad
  \norm{\vc{v}}_p=\Bigl(\sum_{i=1}^{n}\lvert v_i\rvert^{p}\Bigr)^{1/p},
  \qquad
  \norm{\vc{v}}_\infty=\max_{1\le i\le n}\lvert v_i\rvert,
  \qquad
  \cos(\vc{u},\vc{v})=\frac{\ip{\vc{u}}{\vc{v}}}{\norm{\vc{u}}\norm{\vc{v}}} .
\]
Write $\vc{e}_i$ for the standard basis vector with a $1$ in slot $i$, and
$\one=[\,1,\,1,\,\dots,\,1\,]$ for the all-ones vector in $\R^{n}$. In $\R^{n}$
the angle between nonzero vectors is \emph{defined} as
$\theta=\arccos\cos(\vc{u},\vc{v})\in[0,\pi]$ (\S24.5). The sum formulas
$\sum_{i=1}^{n}i=\tfrac{n(n+1)}{2}$ and
$\sum_{i=1}^{n}i^{2}=\tfrac{n(n+1)(2n+1)}{6}$ are in the foundations review.

\medskip
\noindent Cauchy--Schwarz holds in $\R^{n}$ (\S24.3). Your proof of Homework 1
Problem E.2 used only \textbf{D1}--\textbf{D4}, so it already proves it.

%------------------------------------------------------------------------------
\begin{prob}[Arithmetic in $\R^{4}$]
Let
\[
  \vc{u}=[\,1,\,-2,\,0,\,3\,],\qquad
  \vc{v}=[\,2,\,1,\,-1,\,0\,],\qquad
  \vc{w}=[\,0,\,3,\,2,\,-1\,],\qquad
  \alpha=-2 .
\]
\begin{enumerate}[label=\textbf{(\alph*)},leftmargin=2.6em,itemsep=7pt,topsep=6pt]
  \item $\vc{u}+\vc{v}$, \ $\vc{u}-\vc{v}$, \ $\alpha\vc{w}$, \ and
        $2\vc{u}-\vc{v}+3\vc{w}$.
  \item $\norm{\vc{u}}$, $\norm{\vc{v}}$, $\norm{\vc{w}}$, exactly.
  \item $\norm{\vc{u}}_1$, $\norm{\vc{u}}_\infty$, $\norm{\vc{w}}_1$,
        $\norm{\vc{w}}_\infty$.
  \item $\ip{\vc{u}}{\vc{v}}$, $\ip{\vc{u}}{\vc{w}}$, $\ip{\vc{v}}{\vc{w}}$.
        Which pair is orthogonal?
  \item $\cos(\vc{u},\vc{w})$ and $\cos(\vc{v},\vc{w})$, exactly, and the angle
        between $\vc{u}$ and $\vc{w}$ in degrees.
  \item $d_2(\vc{u},\vc{w})$, $d_1(\vc{u},\vc{w})$ and
        $d_\infty(\vc{u},\vc{w})$, and confirm they come out in the order
        \S14.8 predicts.
  \item $\vc{u}$ written as a combination of $\vc{e}_1,\dots,\vc{e}_4$, and
        $\ip{\vc{e}_4}{\vc{u}}$.
\end{enumerate}
\end{prob}

%------------------------------------------------------------------------------
\begin{prob}[The dot-product laws in $\R^{n}$]
\S24.3 says the laws survive in $\R^{n}$ because the proofs ``run over $n$ slots
instead of $2$''. Run them. Use $\Sigma$ notation throughout, and at each step
name the fact you use: a property of real numbers, or one of the rules for sums
in the foundations review (splitting a sum, pulling out a constant).
\begin{enumerate}[label=\textbf{(\alph*)},leftmargin=2.6em,itemsep=7pt,topsep=6pt]
  \item \textbf{D1}: $\ip{\vc{u}}{\vc{v}}=\ip{\vc{v}}{\vc{u}}$.
  \item \textbf{D2}: $\ip{\vc{u}}{\vc{v}+\vc{w}}=\ip{\vc{u}}{\vc{v}}+\ip{\vc{u}}{\vc{w}}$.
  \item \textbf{D3}: $\ip{\vc{u}}{\alpha\vc{v}}=\alpha\ip{\vc{u}}{\vc{v}}$.
  \item \textbf{D4}: $\ip{\vc{u}}{\vc{u}}\ge0$, with equality only for
        $\vc{u}=\vc{0}$.
  \item \textbf{N2} and \textbf{N3} for the $1$-norm on $\R^{n}$:
        $\norm{\alpha\vc{v}}_1=\lvert\alpha\rvert\norm{\vc{v}}_1$ and
        $\norm{\vc{u}+\vc{v}}_1\le\norm{\vc{u}}_1+\norm{\vc{v}}_1$.
\end{enumerate}
\end{prob}

%------------------------------------------------------------------------------
\begin{prob}[Norms of patterned vectors]
Let $n\ge1$. Give each answer as a formula in $n$.
\begin{enumerate}[label=\textbf{(\alph*)},leftmargin=2.6em,itemsep=7pt,topsep=6pt]
  \item $\norm{\one}_1$, $\norm{\one}_2$, $\norm{\one}_\infty$.
  \item Let $\vc{v}=[\,1,\,2,\,3,\,\dots,\,n\,]$. Compute $\norm{\vc{v}}_1$,
        $\norm{\vc{v}}_2$, $\norm{\vc{v}}_\infty$ and $\ip{\one}{\vc{v}}$.
  \item With $\vc{v}$ as in (b), show that
        \[
          \cos(\one,\vc{v}) \;=\; \frac{\sqrt6}{2}\sqrt{\frac{n+1}{2n+1}} .
        \]
        Evaluate the angle between $\one$ and $\vc{v}$ in degrees for $n=1$,
        $2$ and $10$, and find its limit as $n\to\infty$.
  \item For $i\neq j$, compute $\norm{\vc{e}_i-\vc{e}_j}_1$,
        $\norm{\vc{e}_i-\vc{e}_j}_2$ and $\norm{\vc{e}_i-\vc{e}_j}_\infty$.
        Why does the answer not depend on $n$, $i$ or $j$?
\end{enumerate}
\end{prob}

%------------------------------------------------------------------------------
\begin{prob}[Cosine similarity for documents]
A common way to compare documents is to count words. Fix a vocabulary of five
words --- \emph{vector}, \emph{matrix}, \emph{norm}, \emph{pizza},
\emph{soccer} --- and record each document as a vector in $\R^{5}$ of word
counts:
\begin{center}
\renewcommand{\arraystretch}{1.2}
\begin{tabular}{@{}l c@{}}
\toprule
\textbf{Document} & \textbf{Counts} \\
\midrule
$A$ & $[\,3,\,1,\,4,\,0,\,0\,]$ \\
$B$ & $[\,6,\,2,\,8,\,0,\,0\,]$ \\
$C$ & $[\,0,\,0,\,0,\,1,\,5\,]$ \\
$D$ & $[\,2,\,0,\,3,\,1,\,0\,]$ \\
\bottomrule
\end{tabular}
\end{center}
\begin{enumerate}[label=\textbf{(\alph*)},leftmargin=2.6em,itemsep=7pt,topsep=6pt]
  \item Compute $\cos(A,B)$, $\cos(A,C)$, $\cos(A,D)$ and $\cos(C,D)$ exactly,
        and the last two to four decimal places.
  \item Compute the Euclidean distances $d_2(A,B)$, $d_2(A,C)$, $d_2(A,D)$.
  \item Which document is most similar to $A$ by cosine similarity? Which is
        nearest to $A$ by Euclidean distance? Document $B$ is document $A$ pasted
        in twice. Explain which of the two measures gives the sensible answer
        for comparing topics, and why.
\end{enumerate}
\end{prob}

%------------------------------------------------------------------------------
\begin{prob}[The mean and the variance are a projection]
A data set of $n$ numbers is a vector $\vc{x}\in\R^{n}$. Take
$\vc{x}=[\,4,\,7,\,1,\,8,\,5\,]$.
\begin{enumerate}[label=\textbf{(\alph*)},leftmargin=2.6em,itemsep=7pt,topsep=6pt]
  \item Compute $\proj{\one}{\vc{x}}$. Then prove, for any $\vc{x}\in\R^{n}$,
        that $\proj{\one}{\vc{x}}=\bar{x}\,\one$, where
        $\bar{x}=\tfrac1n\sum_{i=1}^{n}x_i$ is the mean. Projecting onto $\one$
        replaces every entry by the average.
  \item The orthogonal piece $\vc{x}-\bar{x}\one$ is the vector of
        \emph{deviations from the mean}. Compute it, check that it is orthogonal
        to $\one$, and deduce the general fact that deviations from the mean sum
        to zero.
  \item The (population) \textbf{variance} is
        $\sigma^{2}=\tfrac1n\norm{\vc{x}-\bar{x}\one}^{2}$ and the
        \textbf{standard deviation} is $\sigma$. Compute both.
  \item For a second data set $\vc{y}$, the \textbf{correlation coefficient}
        of $\vc{x}$ and $\vc{y}$ is the cosine similarity of their deviation
        vectors, $r=\cos(\vc{x}-\bar{x}\one,\ \vc{y}-\bar{y}\one)$. Compute
        $r$ exactly and to four decimal places for
        $\vc{y}=[\,3,\,5,\,1,\,6,\,5\,]$.
\end{enumerate}
\end{prob}

%------------------------------------------------------------------------------
\begin{prob}[The standard basis, and the diagonals of a cube]
\begin{enumerate}[label=\textbf{(\alph*)},leftmargin=2.6em,itemsep=7pt,topsep=6pt]
  \item Show that $\ip{\vc{e}_i}{\vc{v}}=v_i$ for every $\vc{v}\in\R^{n}$, and
        that $\ip{\vc{e}_i}{\vc{e}_j}$ is $1$ when $i=j$ and $0$ otherwise.
  \item Show that the cosine of the angle between the diagonal $\one$ and the
        edge direction $\vc{e}_1$ is $1/\sqrt{n}$. Give the angle in degrees for
        $n=2$, $3$, $4$ and $100$, and find its limit as $n\to\infty$. Say in one
        sentence what this means for the long diagonal of a high-dimensional
        cube.
\end{enumerate}
\end{prob}

%------------------------------------------------------------------------------
\begin{prob}[Comparing the norms in $\R^{n}$]
Prove each inequality for every $\vc{v}\in\R^{n}$. Work from the $\Sigma$
formulas.
\begin{enumerate}[label=\textbf{(\alph*)},leftmargin=2.6em,itemsep=7pt,topsep=6pt]
  \item $\norm{\vc{v}}_\infty\le\norm{\vc{v}}_2\le\norm{\vc{v}}_1$.
        \hint{For the second, compare
        $\norm{\vc{v}}_1^{2}=\bigl(\sum\lvert v_i\rvert\bigr)^{2}$ with
        $\sum v_i^{2}$. What do the cross terms contribute?}
  \item $\norm{\vc{v}}_1\le n\,\norm{\vc{v}}_\infty$ and
        $\norm{\vc{v}}_2\le\sqrt{n}\,\norm{\vc{v}}_\infty$.
  \item $\norm{\vc{v}}_1\le\sqrt{n}\,\norm{\vc{v}}_2$.
        \hint{Apply Cauchy--Schwarz in $\R^{n}$ to the vectors
        $[\,\lvert v_1\rvert,\,\dots,\,\lvert v_n\rvert\,]$ and $\one$.}
\end{enumerate}
\noindent In one sentence: together these say that in $\R^{n}$ any two of the
three norms are within a factor of $n$ of each other. What changes as $n$ grows?
\end{prob}

%==============================================================================
\section{Derivatives, partial derivatives, and the gradient}
%==============================================================================
\emph{Vector Calculus Unit, Part 1.} In this part, a bare section number such as
\S6.1 refers to the \emph{vector calculus} unit, whose numbering restarts at
$1$. References into the vector unit say so. Throughout,
\[
  f'(x)=\lim_{h\to0}\frac{f(x+h)-f(x)}{h},
  \qquad
  \pd{f}{x}(x,y)=\lim_{h\to0}\frac{f(x+h,y)-f(x,y)}{h},
  \qquad
  \grad f=\Bigl[\,\pd{f}{x},\ \pd{f}{y}\,\Bigr],
\]
and the exponential and logarithm laws are \textbf{E1}--\textbf{E5} and
\textbf{L1}--\textbf{L5} of \S3.

%------------------------------------------------------------------------------
\begin{prob}[The rules, in one variable]
Differentiate each function, and name the rule or rules you use --- those of
\S2, and the derivatives of $e^{x}$ and $\ln x$ found in \S4. Simplify where it
helps.
\begin{enumerate}[label=\textbf{(\alph*)},leftmargin=2.6em,itemsep=7pt,topsep=6pt]
  \item $x^{3}e^{x}$
  \item $\ln(x^{2}+1)$
  \item $e^{-x^{2}/2}$
  \item $\dfrac{\sin x}{1+e^{x}}$
  \item $\sqrt{x}\,\ln x$
  \item $\ln(\ln x)$
  \item $e^{3x}\tan x$
  \item $\sec(x^{2})$
\end{enumerate}
\end{prob}

%------------------------------------------------------------------------------
\begin{prob}[Exponential and logarithm algebra]
\begin{enumerate}[label=\textbf{(\alph*)},leftmargin=2.6em,itemsep=7pt,topsep=6pt]
  \item Simplify each expression, naming the law of \S3 used at each step:
        \[
          \ln\bigl(e^{3}e^{-5}\bigr),\qquad
          e^{2\ln3},\qquad
          \ln\sqrt{e},\qquad
          \ln8-3\ln2,\qquad
          e^{\ln5-\ln2}.
        \]
  \item Solve $e^{2x}-3e^{x}+2=0$.
        \hint{It is a quadratic in the unknown $e^{x}$.}
  \item Solve $\ln x+\ln(x-3)=\ln4$. One solution of the quadratic you reach
        must be rejected. Which one, and why?
\end{enumerate}
\end{prob}

%------------------------------------------------------------------------------
\begin{prob}[Two ways, again]
\S4.2 and \S4.3 found the derivative of $e^{x}$ two ways: by implicit
differentiation, and from the definition. Do the same kind of thing here.
\begin{enumerate}[label=\textbf{(\alph*)},leftmargin=2.6em,itemsep=7pt,topsep=6pt]
  \item Let $a>0$. Derive $\dfrac{d}{dx}a^{x}=a^{x}\ln a$ in two ways:
        \begin{enumerate}[label=(\roman*),leftmargin=2.2em,itemsep=3pt,topsep=3pt]
          \item write $a^{x}=e^{x\ln a}$, justifying it by the inverse
                relation $e^{\ln a}=a$ and \textbf{E4}, then use the chain rule;
          \item put $y=a^{x}$, take $\ln$ of both sides using \textbf{L4}, and
                differentiate implicitly.
        \end{enumerate}
  \item \S4.2 went from $\ln$ to $\exp$. Go the other way: taking
        $\tfrac{d}{dx}e^{x}=e^{x}$ as known, derive
        $\tfrac{d}{dx}\ln x=1/x$ by writing $y=\ln x$ as $e^{y}=x$ and
        differentiating implicitly.
\end{enumerate}
\end{prob}

%------------------------------------------------------------------------------
\begin{prob}[Partial derivatives]
Compute $\partial f/\partial x$ and $\partial f/\partial y$. Assume each
function is considered where it is defined.
\begin{enumerate}[label=\textbf{(\alph*)},leftmargin=2.6em,itemsep=7pt,topsep=6pt]
  \item $f(x,y)=3x^{2}y^{3}-4xy+7$
  \item $f(x,y)=x\,e^{xy}$
  \item $f(x,y)=\ln(x^{2}+y^{2})$
  \item $f(x,y)=\sin x\cos y$
  \item $f(x,y)=x/y$
  \item $f(x,y)=(x-2y)^{5}$
  \item $f(x,y)=e^{x}\ln y$
  \item $f(x,y)=\sqrt{x^{2}+y^{2}}$
\end{enumerate}

\heads{In (b), the variable held constant appears in two places --- in front, and
inside the exponent. Holding $y$ fixed, $f$ is a product of two functions of
$x$, so the product rule is needed. Holding $x$ fixed, the front $x$ is a
constant.}
\end{prob}

%------------------------------------------------------------------------------
\begin{prob}[Partial derivatives in context]
\begin{enumerate}[label=\textbf{(\alph*)},leftmargin=2.6em,itemsep=7pt,topsep=6pt]
  \item The volume of a cylinder of radius $r$ and height $h$ is
        $V(r,h)=\pi r^{2}h$. Compute $\partial V/\partial r$ and
        $\partial V/\partial h$, evaluate both at $r=3$~cm, $h=10$~cm, and
        give their units. Use them to estimate which adds more volume: widening
        the radius by $1$~mm or raising the height by $1$~mm. Then compute the
        two actual changes and compare.
  \item For a fixed amount of gas, pressure depends on temperature $T$ (in
        kelvin) and volume $V$ (in litres) as $P(T,V)=8T/V$ (in kilopascals).
        Compute $P$, $\partial P/\partial T$ and $\partial P/\partial V$ at
        $T=300$, $V=10$, with units. Explain the sign of each partial derivative
        physically.
\end{enumerate}
\end{prob}

%------------------------------------------------------------------------------
\begin{prob}[The gradient]
\begin{enumerate}[label=\textbf{(\alph*)},leftmargin=2.6em,itemsep=7pt,topsep=6pt]
  \item Let $f(x,y)=x^{3}-3xy+y^{2}$. Compute $\grad f$, evaluate it at $(1,2)$,
        $(2,-1)$ and $(0,0)$, and compute $\norm{\grad f(2,-1)}$.
  \item Find every point where $\grad f=\vc{0}$ for the $f$ of (a).
  \item Let $f(x,y)=\ln(x^{2}+y^{2})$. Show that
        $\norm{\grad f(x,y)}=2/\norm{[x,y]}$, and describe the direction of
        $\grad f$.
  \item Let $f(x,y)=\sqrt{x^{2}+y^{2}}=\norm{[x,y]}$. Show that
        $\grad f(x,y)$ is the \emph{normalization} of $[\,x,\,y\,]$ (vector unit
        \S12.1), so that $\norm{\grad f}=1$ at every point other than the
        origin. Why is $\grad f$ undefined at the origin?
\end{enumerate}
\end{prob}

\bigskip
\noindent\rule{\textwidth}{0.4pt}

\smallskip
\noindent{\small Questions about any problem here belong in the class Discord,
where everyone can see the answer. Questions that are specific to your own work
--- ``is my proof of C.6(b) valid?'' --- are better brought to office hours,
Thursdays 3:00--4:00 PM in Kiely Hall 508. Please email ahead.}

\end{document}
